% !TeX program = tectonic
\documentclass[11pt,a4paper]{article}
\usepackage{fontspec}
\usepackage[russian]{babel}
\babelfont{rm}[Language=Default]{Liberation Serif}
\babelfont[Language=Default]{sf}{Carlito}
\babelfont[Language=Default]{tt}{DejaVu Sans Mono}
\usepackage{amsmath,amssymb,amsthm}
\usepackage{geometry}
\geometry{a4paper, top=2.4cm, bottom=2.4cm, left=2.4cm, right=2.4cm}
\usepackage{booktabs,tabularx,enumitem,xcolor,tcolorbox}
\tcbuselibrary{breakable,skins}
\definecolor{linkblue}{HTML}{1F4E79}
\definecolor{statgreen}{HTML}{2F6B3A}
\definecolor{goldacc}{HTML}{8B7E5A}
\usepackage[colorlinks=true,linkcolor=linkblue,urlcolor=linkblue,unicode=true]{hyperref}
\theoremstyle{plain}
\newtheorem{theorem}{Теорема}
\newtheorem{lemma}[theorem]{Лемма}
\newtheorem{proposition}[theorem]{Предложение}
\newtheorem{corollary}[theorem]{Следствие}
\theoremstyle{definition}
\newtheorem{definition}[theorem]{Определение}
\theoremstyle{remark}
\newtheorem{remark}[theorem]{Замечание}
\newtcolorbox{statusbox}[1][]{colback=statgreen!5,colframe=statgreen!75,
  title={\textbf{Сводка доказанного:} #1},fonttitle=\small,boxrule=0.7pt,arc=2pt,breakable}
\newtcolorbox{databox}[1][]{colback=goldacc!6,colframe=goldacc!80,
  title={\textbf{Данные протокола:} #1},fonttitle=\small,boxrule=0.7pt,arc=2pt,breakable}
\newtcolorbox{verifybox}[1][]{colback=linkblue!5,colframe=linkblue!80,
  title={\textbf{Верификация:} #1},fonttitle=\small,boxrule=0.7pt,arc=2pt,breakable}
\widowpenalty=10000
\clubpenalty=10000
\setlength{\parskip}{0.35em}
\newcommand{\CC}{\mathbb{C}}
\newcommand{\RR}{\mathbb{R}}
\newcommand{\ZZ}{\mathbb{Z}}
\newcommand{\QQ}{\mathbb{Q}}
\newcommand{\Ga}{\Gamma}
\newcommand{\Bfun}{\mathrm{B}}
\newcommand{\im}{\operatorname{Im}}
\newcommand{\re}{\operatorname{Re}}
\newcommand{\lcm}{\operatorname{lcm}}
\newcommand{\SNF}{\operatorname{SNF}}
\newcommand{\Jac}{\operatorname{Jac}}
\newcommand{\rank}{\operatorname{rank}}
\newcommand{\Ind}{\operatorname{Index}}
\newcommand{\disc}{\operatorname{disc}}
\begin{document}
\begin{center}
{\LARGE\bfseries Сертификат G: индексы решётки Ходжа}\\[0.4em]
{\large спектр решёток, CM-дроби и слой Коула–Селберга}\\[0.6em]
{\normalsize Исаев Исхак Хамзатович}\\[0.2em]
{\small Программа <<Динамический принцип>> --- лаборатория \texttt{hodge-laboratory}}\\[0.15em]
{\small Отдельная монография теоремы · Монография 15/16}
\end{center}
\vspace{0.4em}
{\color{linkblue}\hrule height 0.8pt}
\vspace{0.8em}

\section{Постановка}

Сертификат G даёт универсальную нормировку спектра произвольной
решётки и вычисляет индексы решётки Ходжа в числовых полях ---
мост от целочисленной геометрии к $\Ga$-нормировкам сертификата H.

\begin{theorem}[G1--G6]\label{th:Gstand}
\begin{enumerate}[leftmargin=2em]
\item[(G1)] Универсальная нормировка спектра:
$\lambda_{m,n}=(2\pi)^2|m\tau-n|^2/(\im\tau)^2$; множитель $4\pi^2$
--- универсальный Tate-фактор.
\item[(G2)] Для приведённых решёток $\lambda_1=4\pi^2/(\im\tau)^2$;
для CM-полей --- точная дробь: семейства $16\pi^2/d$ и $4\pi^2/d$.
Уровни $\pi/15$ и $\pi/30$ --- спектральные единицы решёток
$\ZZ[\sqrt{-15}]$, $\ZZ[\sqrt{-30}]$: $\lambda_1/(4\pi)=\pi/N$.
\item[(G3)] Индексы в числовых полях восстанавливаются с
единственностью $2\varepsilon<1/Q$ (поле-LLL с тройкой $(p,q,s)$).
\item[(G4)] Слой Коула--Селберга: $\omega(i)^2=\Ga(1/4)^4/(16\pi)$;
$\omega(\tau_7)^2=\frac{21+\sqrt{-7}}{896}
\bigl[\Ga(\frac17)\Ga(\frac27)\Ga(\frac47)\bigr]^2/\pi^2$;
$\omega(\tau_3)^2=\frac{3-\sqrt{-3}}8
\Ga(\frac13)^6/(\pi^2 2^{8/3})$; для $d=15$: $j$-пара в
$\QQ(\sqrt5)$ и $R=\frac{(2\sqrt5-3)+(\sqrt5-2)\sqrt{-3}}2$.
\item[(G5)] K3-мост: квадрант-интеграл $\Ga(1/4)^4/(16\pi\sqrt2)$;
индекс $\sqrt2/32$ в $\Ga$-шкале.
\item[(G6)] Лестница $\mu_N$-квантов: квант $2\pi/N$, спиновая
половина $\pi/N$; $\delta_C=\pi/7$ --- ступень $N=7$; $\pi/15$,
$\pi/30$ --- следующие ступени.
\end{enumerate}
\end{theorem}

\begin{proof}
(G1) Плоский лапласиан на эллиптической кривой с решёткой
$\ZZ\tau+\ZZ$: собственные значения $\propto|m\tau-n|^2$;
нормировка на $(2\pi)^2$ даёт формулу.
(G2) Для приведённого $\tau$ минимум $|m\tau-n|^2=1$ достигается на
кратчайшем векторе; CM-случай: $\tau=\sqrt{-d}$ либо
$\frac{1+\sqrt{-d}}2$, дробь раскрывается в точное $4\pi^2/d$ или
$16\pi^2/d$ по классификации кратчайших векторов; для $d=15$
прямое вычисление нормы даёт $\lambda_1/(4\pi)=\pi/15$.
(G3) Разделение индексов в поле оценивается снизу через
дискриминант; LLL в базисе $(p,q,s)$ уникальна при
$2\varepsilon<1/Q$; остаток восстановления $2.31\cdot10^{-102}$.
(G4) Формулы Коула--Селберга: квадрат периода --- произведение
$\Ga$-значений с рациональными показателями; все тождества
проверены с остатками $<10^{-100}$.
(G5) Разбиение интеграла на квадранты и замена переменных дают
точное значение; остаток $3.92\cdot10^{-43}$.
(G6) Спиновая половина --- следствие двойного накрытия; лестница
согласована с тройками стендов.
\end{proof}

\begin{databox}[числа]
Не-CM тест: минимальный многочлен $4X^3-1$ --- поле-LLL отличает
CM от не-CM. Поле-LLL $d=15$: восстановление
$\frac{3+5\sqrt{-15}}8$. Спектральные дроби подтверждены на девяти
решётках G1; семейства $16\pi^2/d$, $4\pi^2/d$ --- точные.
\end{databox}

\begin{statusbox}[итог]
Сертификат G принят ($6/6$): универсальная нормировка, точные
дробности, поле-LLL, слой Коула--Селберга, K3-мост, лестница ---
все пункты сертифицированы.
\end{statusbox}

\begin{verifybox}[как воспроизвести]
\texttt{python3 laboratory.py} $\to$ <<Сертификаты $\to$ G>>;
конструктор: произвольный $d$ --- лаборатория вычисляет
$\lambda_1$ и печатает дробь.
\end{verifybox}

\vfill
{\color{linkblue}\hrule height 0.6pt}
\vspace{0.5em}
{\small\color{gray}
Лицензия: индивидуальная исключительная лицензия Исаева Исхака Хамзатовича.
Все права на вычисления, формулы и тексты принадлежат автору.
Верификация: \texttt{python3 laboratory.py} (пункт <<Стенды>>/<<Сертификаты>>),
Lean: \texttt{verification/lean/HodgeLaboratory.lean}.
}
\end{document}
